% Patch record for R2b-04 -- Hardware parameter subsection
% Reviewer #2 (round 2). Applied to sections/methodology.tex, 2026-10-27.
% Target: \subsection{Physical Platform Construction and Experimental Setup} (ssec:Construction).
% Changes: (1) 'main control and processing unit' -> 'onboard control unit' + sentence stating the
%   neural circuits run on an external computer; (2) 3D-printing sentence now states all parts are
%   PLA (previously PLA/PETG split); (3) new \subsubsection{Hardware Specification}
%   (ssec:HardwareSpec) with table:HardwareParams, inserted after fig:robot_real.
% Data source: author-confirmed 2026-10-27 -- total mass 1.3 kg (battery incl.), all-PLA structure,
%   Logitech C270, network on external laptop, UBEC included, Hall sensor excluded. Hardware
%   specs from thesis T11498 (Table 3, Table 8), read-only reference.
% Unified diff below (context-free):
% diff --git i/sections/methodology.tex w/sections/methodology.tex
% index 7093ddf..41602e0 100644
% --- i/sections/methodology.tex
% +++ w/sections/methodology.tex
% @@ -474 +474 @@ Based on the design presented in subsection~\ref{ssec:Platform_Design} and in Fi
% -The ESP32-S3 Zero microcontroller was used as the main control and processing unit, responsible for sensor acquisition and actuator control. The latter was intermediated by the PCA9685 servo controller and the mini-L298N module, which regulated the speed of the N20 DC motors. 
% +The ESP32-S3 Zero microcontroller was used as the \revblue{onboard control unit}, responsible for sensor acquisition and actuator control. The latter was intermediated by the PCA9685 servo controller and the mini-L298N module, which regulated the speed of the N20 DC motors. \revblue{The neural circuits themselves run on an external computer, as detailed in Section~\ref{ssec:HardwareSpec}.} 
% @@ -478 +478 @@ The network inputs were implemented through an MPU-6050 inertial sensor, an HMC5
% -The structural components of the limbs and chassis were manufactured using 3D printing. Polylactic acid (PLA), filament was employed for parts not subjected to significant mechanical loads, while polyethylene terephthalate glycol-modified was used for the limbs and the central chassis section exposed to higher stresses during locomotion. Additionally, a rubber stopper was incorporated at the end of each limb, at the ground contact point, to increase traction during quadrupedal gait. The fully assembled robot is shown in Figure~\ref{fig:robot_real}. 
% +\revblue{All structural components of the limbs and chassis were manufactured by 3D printing in polylactic acid (PLA).} Additionally, a rubber stopper was incorporated at the end of each limb, at the ground contact point, to increase traction during quadrupedal gait. The fully assembled robot is shown in Figure~\ref{fig:robot_real}. 
% @@ -486,0 +487,25 @@ The structural components of the limbs and chassis were manufactured using 3D pr
% +\subsubsection{\revblue{Hardware Specification}}
% +\label{ssec:HardwareSpec}
% +
% +\revblue{Table~\ref{table:HardwareParams} consolidates the main parameters of the physical platform. The ESP32-S3 Zero is the only computing element mounted on the robot: it acquires the inertial data, drives the actuators, and exchanges commands and sensor readings with an external computer over a serial link. The obstacle-sensory, basal-ganglia, locomotion, and gait-decision circuits, together with the MLP terrain classifier and the image processing, are executed on that external laptop (Intel Core i7-1065G7, 12~GB RAM, Linux), not on the robot itself.}
% +
% +\begin{table}[t]
% +\centering
% +\color{blue}
% +\caption{Physical platform parameters.}
% +\begin{tabularx}{\textwidth}{|l|>{\raggedright\arraybackslash}X|}
% +\hline
% +\textbf{Parameter} & \textbf{Value} \\ \hline
% +Total mass (battery included) & 1.3~kg \\ \hline
% +Structural material & PLA, 3D printed (all parts) \\ \hline
% +Joint actuators & 12 MG996R servomotors (3 per leg), $11~\text{kg·cm}$ nominal torque, driven by a PCA9685 controller \\ \hline
% +Wheel actuators & 4 N20 geared DC motors, $104$~RPM (1 per leg), driven by a mini-L298N module \\ \hline
% +Onboard computing platform & ESP32-S3 Zero (dual-core, 240~MHz, 2~MB PSRAM, 4~MB flash) \\ \hline
% +Off-board computing platform & Laptop, Intel Core i7-1065G7, 12~GB RAM; serial link to the robot \\ \hline
% +Sensors & MPU-6050 IMU, HMC5883L magnetometer, Logitech C270 camera \\ \hline
% +Battery & Li-ion, 7.2~V, 3300~mAh \\ \hline
% +Voltage regulation & 8~A switching UBEC, 5~V output \\ \hline
% +\end{tabularx}
% +\label{table:HardwareParams}
% +\end{table}
% +
%
% Follow-up edit (same row), sections/results.tex, ~line 737 (MLP classifier-comparison paragraph):
%  - 'Given the hard real-time constraints of the neural-arbitration pipeline on the resource-
%    constrained embedded microcontroller' -> 'Given the real-time constraints ... and the prospect of
%    deploying the classifier on a resource-constrained embedded microcontroller' (the ESP32-S3 is a
%    deployment target evaluated on-device, not where the network ran in the experiments).
%  - 'host PC (PyTorch)' now points to ssec:HardwareSpec (same laptop as the physical experiments).
